\begin{enumerate}
	\item 	\begin{itemize}[label =$\bullet$]
				\item Le côté le plus long est [BC] et $\text{BC}^2=8^2=64$.
				\item D'autre part : $\text{AC}^2+\text{AB}^2=4,8^2+6,4^2=64$.
		\end{itemize}
		Donc $\text{BC}^2=\text{AC}^2+\text{AB}^2$.

		D'après la réciproque du théorème de Pythagore, on conclut que le triangle ABC est rectangle en A.
		\item On sait que :
		\begin{itemize}[label =$\bullet$]
			\item Les points C, A et E d'une part et B, A et D d'autre part, sont alignés dans le même ordre.
				\item (BC) et (DE) sont parallèles;
		\end{itemize}
		D'après le théorème de Thalès, on a :

		\(\begin{array}{ccccc}
		\dfrac{\text{CA}}{\text{AE}}&=&\dfrac{\text{BA}}{\text{AD}}&=&\dfrac{\text{BC}}{\text{DE}}\\
		 & & & & \\
		\dfrac{4,8}{\text{AE}}&=&\dfrac{6,4}{4,8}&=&\dfrac{8}{\text{DE}}\\
		\end{array}\)

		Donc $\text{DE}=\dfrac{8\times 4,8}{6,4}=6$ cm,

		et $\text{AE}=\dfrac{4,8\times 4,8}{6,4}=3,6$ cm.
		\item Les droites (BC) et (DE) sont parallèles et sont coupées par une sécante (DB), donc les angles alternes-internes formés sont égaux. Ainsi $\widehat{\text{ABC}}=\widehat{\text{ADE}}$.
		\item On sait que $\widehat{\text{ABC}}=\widehat{\text{ADE}}$ et que $\widehat{\text{BAC}}=\widehat{\text{DAE}}$ (ce sont des angles droits).

		Donc $\widehat{\text{ACB}}=\widehat{\text{AED}}$.

		Ainsi les triangles ABC et ADE ont leurs angles deux à deux de même mesure, les triangles sont donc semblables.
		\item
		L'aire du quadrilatère est la somme des aires des quatre triangles.

		 $\mathcal{A}_{\text{BCDE}}=\mathcal{A}_{\text{ABC}}+\mathcal{A}_{\text{ABE}}+\mathcal{A}_{\text{ADE}}+\mathcal{A}_{\text{ADC}}$.

		 Ce sont des triangles rectangles, donc l'aire est la moitié du produit des longueurs des côtés de l'angle droit.
		\begin{itemize}[label =$\bullet$]
			\item $\mathcal{A}_{\text{ABC}}=\dfrac{4,8\times 6,4}{2}=15,36$ cm$^2$.
			\item $\mathcal{A}_{\text{ABE}}=\dfrac{3,6\times 6,4}{2}=11,52$ cm$^2$.
			\item $\mathcal{A}_{\text{ADE}}=\dfrac{4,8\times 3,6}{2}=8,64$ cm$^2$.
			\item $\mathcal{A}_{\text{ADC}}=\dfrac{4,8\times 4,8}{2}=11,52$ cm$^2$.
		\end{itemize}
		Ainsi $\mathcal{A}_{\text{BCDE}}=15,36+11,52+8,64+11,52=47,04$ cm$^2$.
\end{enumerate}
\needspace{10\baselineskip}
% saut de page si moins de 10 lignes disponibles
